% Introduction section. Paste into your own thesis, or \input it.
%
% STRUCTURE (after Aida Jimenez's MSc thesis ch. 1, Former_students_work/): her chapter runs
% funnel -> state of the art -> "Aim of the Thesis". Ours splits the same way, except the state of
% the art is its own chapter because objective 4 demands more than three pages of it. This section
% is her 1.2: the aim in prose, then what will actually be found out, then what bounds it.
%
% What her aim section does NOT have, and this one does: research questions with an order and a
% stated dependency, and an explicit limits paragraph. What it has that an earlier draft of ours
% lacked: a plain statement of the aim before any of the detail.
%
% TODO (CGPS): no exact figure from the Copenhagen General Population Study appears here --
% cohort size, repeat-scan counts, follow-up. Data access is pending.

\section{Aim of the thesis}
\label{sec:aims}

The aim of this thesis is to describe the topology of the coronary artery tree across a population,
and to establish whether an individual tree can be called unusual against that description.
Three things are needed together and do not yet exist that way: a construction that turns a
segmented coronary tree into a measurable object, features on that object that a machine computes
identically every time, and a cohort large enough and healthy enough for the resulting distribution
to carry meaning.

The work runs in two stages.
ImageCAS-X supplies 800 annotated cases on which the construction and the features are built, and
on which their reproducibility is measured.
The Copenhagen General Population Study then supplies what ImageCAS-X cannot: a cohort not selected
for disease, at a scale large enough to describe a distribution rather than a sample, and
adjudicated outcomes against which any extreme value in that distribution can be tested
\cite{fuchs2023cgps}.

Four questions run in order, and each is worth asking only once the one before it has an answer.

The first is what can be measured at all.
A coronary tree becomes a measurable object once it is a graph: vessels as edges, bifurcations and
terminations as nodes, rooted at the ostium where the artery leaves the aorta.
That construction is exact rather than approximate in ImageCAS-X, whose centerlines meet by sharing
a point index and so carry the connectivity without any distance threshold
\cite{bransby2026imagecasx}.
What such a graph supports is still open.
Branching pattern, bifurcation angle, tortuosity and dominance each need one definition that a
machine computes identically every time.
Some of them may not survive that requirement.

The second is how well those features can be measured, and the answer decides what the rest is
worth.
Each feature carries two error terms, and neither appears in the coronary literature found by the
searches recorded for this thesis.
The first is segmentation error, since the tree will be extracted from a predicted lumen rather than
a drawn one.
Measuring it means computing every feature on both and comparing them over the 160 test cases the
benchmark holds back.
The second is repeated measurement: scanning the same person twice should return the same tree.
Whatever difference remains is the floor beneath which no real change can be detected.
The closest published equivalent is carotid rather than coronary, 177 subjects imaged twice a mean of
$130.2 \pm 8.1$ months apart, whose bifurcation angles had widened significantly by the second scan
\cite{kwon2022carotid}.
The repeat scans in the Copenhagen General Population Study make the coronary version possible.
Until that floor is known, a wide bifurcation and a measurement error are the same number.

The third is what a population of coronary trees looks like.
Normal coronary anatomy has been described before, but in selected samples of a few hundred: over 300
angiograms chosen for absent calcium and no stenosis \cite{medranogracia2016atlas}, and 281 patients
characterized automatically \cite{nannini2024characterization}.
Neither covers the features extracted here, at the scale needed to call a value unusual.
So the distributions have to be described first: how branching, angle and tortuosity are spread, and
how much of that spread belongs to dominance, sex and age rather than to disease.
Dominance is the test of whether the description can be trusted, because ImageCAS-X records it twice
by independent routes \cite{bransby2026imagecasx}.
It is annotated explicitly, and it is implied by which arteries carry labels at all.

The fourth is whether the tails of those distributions mean anything.
An extreme value is only interesting if it exceeds the measurement floor established by the second
question.
It is only useful if it separates people who go on to have events from people who do not.
The Copenhagen cohort carries adjudicated outcomes, so the last step asks whether topological
extremes add anything to the risk factors already used in clinic \cite{fuchs2023cgps}.
A negative answer is a result, provided the second question has been answered well enough to make
the negative believable.

Two limits bound every answer above.
The first is measurement.
Phantoms fabricated with known bifurcation angles were recovered to within a few degrees, but the
same two accepted measurement techniques applied to 50 real scans disagreed with each other by
$12.0^\circ \pm 10.6^\circ$ \cite{givehchi2018phantom}, almost the entire width of the clinical
gradient in Section~\ref{sec:framing}, $12^\circ$ from a normal artery to a significant stenosis.
Every feature in this thesis is therefore computed by one automatic definition applied to the
centerline, and its reproducibility is measured rather than assumed.
The second returns to remodeling.
Because a tree's shape keeps changing as disease grows inside it (Section~\ref{sec:framing}), a
healthy tree found unusual today is not guaranteed to still be unusual, or still healthy, years
later.
That is why the outcomes the Copenhagen cohort carries, and not a single snapshot, are what will
decide whether an unusual shape means anything.
